%% kapa: exam=1 round=30 year=2019 period=1 form=A subject=real_estate subject_ko=부동산학원론 seq=7 number=87
%% source: Q-Net 감정평가사 기출문제 (cst00309) artlSeq=5205205 file=extracted/1008/5205205/2019년 제30회 감정평가사 1차시험 문제지/2019년 제30회 감정평가사 1차시험 문제지/2019년도 제30회 감정평가사 1차 1교시 A형.hwp
%% conversion: kordoc 4.24.51 -> tools/gen1.py (auto); review=proofread-sonnet; flags=none
\begin{question}{87}
부동산투자에서 레버리지(leverage)에 관한 설명으로 옳지 않은 것은?
\begin{choices}
  \item 총투자수익률에서 지분투자수익률을 차감하여 정(+)의 수익률이 나오는 경우에는 정(+)의 레버리지가 발생한다.
  \item 차입이자율이 총투자수익률보다 높은 경우에는 부(-)의 레버리지가 발생한다.
  \item 정(+)의 레버리지는 이자율의 변화 등에 따라 부(-)의 레버리지로 변화될 수 있다.
  \item 부채비율이 상승할수록 레버리지 효과로 인한 지분투자자의 수익률 증대효과가 있지만, 한편으로는 차입금리의 상승으로 지분투자자의 수익률 감소효과도 발생한다.
  \item 대출기간 연장을 통하여 기간이자 상환액을 줄이는 것은 부(-)의 레버리지 발생시 적용할 수 있는 대안 중 하나이다.
\end{choices}
\end{question}
\answer{1}
